% ============================================================
% FIG 6 — STBV-Bench Generation Pipeline
% Shows how VeReMi records flow through the semantic
% transformation engine to produce the benchmark.
% All numbers from Section V of the paper.
% ============================================================
% PREAMBLE:
%   \usepackage{tikz}
%   \usepackage{xcolor}
%   \usetikzlibrary{arrows.meta, positioning, shapes.geometric, calc}
%
% PLACEMENT: full width
%   \begin{figure*}[t]
%     \centering
%     % ============================================================
% FIG 6 — STBV-Bench Generation Pipeline
% Shows how VeReMi records flow through the semantic
% transformation engine to produce the benchmark.
% All numbers from Section V of the paper.
% ============================================================
% PREAMBLE:
%   \usepackage{tikz}
%   \usepackage{xcolor}
%   \usetikzlibrary{arrows.meta, positioning, shapes.geometric, calc}
%
% PLACEMENT: full width
%   \begin{figure*}[t]
%     \centering
%     % ============================================================
% FIG 6 — STBV-Bench Generation Pipeline
% Shows how VeReMi records flow through the semantic
% transformation engine to produce the benchmark.
% All numbers from Section V of the paper.
% ============================================================
% PREAMBLE:
%   \usepackage{tikz}
%   \usepackage{xcolor}
%   \usetikzlibrary{arrows.meta, positioning, shapes.geometric, calc}
%
% PLACEMENT: full width
%   \begin{figure*}[t]
%     \centering
%     % ============================================================
% FIG 6 — STBV-Bench Generation Pipeline
% Shows how VeReMi records flow through the semantic
% transformation engine to produce the benchmark.
% All numbers from Section V of the paper.
% ============================================================
% PREAMBLE:
%   \usepackage{tikz}
%   \usepackage{xcolor}
%   \usetikzlibrary{arrows.meta, positioning, shapes.geometric, calc}
%
% PLACEMENT: full width
%   \begin{figure*}[t]
%     \centering
%     \input{fig6_stbvbench_pipeline.tex}
%     \caption{STBV-Bench generation pipeline. Real VeReMi Extension
%     trajectories supply kinematics; the seeded semantic transformation
%     engine populates free-text scene-context fields across 20 attack
%     families plus a benign control. Kinematic ground truth is preserved
%     separately and never merged with semantic labels.}
%     \label{fig:stbvbench_pipeline}
%   \end{figure*}
% ============================================================

\begin{tikzpicture}[
  arrow/.style={-{Stealth[length=5pt]}, thick, gray!55},
  % [LAYOUT] Stage 1 (Dataset) and Stage 6 (STBV-Bench) blocks kept at
  %          their original widths (2.8cm and 3.0cm respectively).
  %          Intermediate stages reduced from 2.8cm to 2.3cm to recover
  %          horizontal space without altering the visual hierarchy.
  %          The Engine remains the widest and tallest block.
  mainblock/.style={
    rectangle, rounded corners=3pt,
    minimum width=1.85cm, minimum height=1.0cm,
    font=\scriptsize\bfseries, align=center,
    draw, thick
  },
  subblock/.style={
    rectangle, rounded corners=2pt,
    minimum width=1.8cm, minimum height=0.48cm,
    font=\tiny, align=center, text width=1.9cm,
    draw=gray!50, fill=gray!8, thin
  },
  statnode/.style={
    font=\tiny\bfseries, align=center, text width=2.4cm
  },
  % [LAYOUT] Reduced inter-block spacing from 0.55cm to 0.35cm.
  node distance=0.5cm and 0.35cm
]

\definecolor{veremiblue}{RGB}{60,100,180}
\definecolor{engineorange}{RGB}{200,100,20}
\definecolor{validgreen}{RGB}{40,140,60}
\definecolor{benchpurple}{RGB}{120,40,160}
\definecolor{evalcyan}{RGB}{20,120,140}

% ============================================================
% STAGE 1: VeReMi Source
% [LAYOUT] Dataset block kept at original 2.8cm width (unchanged)
% ============================================================
\node[mainblock,
  fill=veremiblue!12, draw=veremiblue, text=veremiblue,
  minimum width=2.8cm, minimum height=1.2cm
] (veremi) {VeReMi Extension\\Dataset};

\node[statnode, below=0.10cm of veremi] (veremi_stat)
  {\textcolor{veremiblue}{221,125 records}\\
   \textcolor{black}{360 vehicles}\\
   \textcolor{black}{180 attackers}};

% ============================================================
% STAGE 2: Field Extraction
% ============================================================
\node[mainblock,
  right=0.30cm of veremi,
  fill=gray!8, draw=gray!50,
  minimum height=1.2cm
] (extract) {Field\\Extraction};

\node[subblock, below=0.08cm of extract] (extract_sub)
  {pos, speed, heading,\\timestamp\\(unmodified)};

% ============================================================
% STAGE 3: Semantic Transformation Engine
% [LAYOUT] Engine kept as largest/widest block at 3.2cm wide,
%          2.2cm tall, to preserve its focal role in the figure
% ============================================================
\node[mainblock,
  right=0.30cm of extract,
  fill=engineorange!10, draw=engineorange, text=engineorange,
  minimum height=3.2cm, minimum width=3.8cm
] (engine) {};

% Explicitly position every title line to guarantee no overlaps
\node[font=\scriptsize\bfseries, text=engineorange, align=center,
      at=(engine.center), yshift=1.0cm]
  {Semantic};
\node[font=\scriptsize\bfseries, text=engineorange, align=center,
      at=(engine.center), yshift=0.65cm]
  {Transformation};
\node[font=\scriptsize\bfseries, text=engineorange, align=center,
      at=(engine.center), yshift=0.30cm]
  {Engine};

% Sub-labels inside engine — positioned explicitly below the title
\node[font=\tiny, text=engineorange!80, align=center,
      at=(engine.center), yshift=-0.25cm]
  {seed = fixed};
\node[font=\tiny, text=engineorange!80, align=center,
      at=(engine.center), yshift=-0.65cm]
  {rule-based, templated};
\node[font=\tiny, text=engineorange!80, align=center,
      at=(engine.center), yshift=-1.05cm]
  {no LLM, no paraphrase};

% Attack families annotation (above engine)
\node[
  font=\tiny\bfseries, text=engineorange,
  align=center, text width=3.2cm,
  above=0.15cm of engine
] {20 attack families + 1 benign control = 21 transformation rules};

% ============================================================
% STAGE 4: Structural Validation
% ============================================================
\node[mainblock,
  right=0.30cm of engine,
  fill=validgreen!10, draw=validgreen, text=validgreen,
  minimum height=1.2cm
] (valid) {Structural\\Validation};

\node[subblock, below=0.08cm of valid] (valid_sub)
  {schema check\\abort on failure};

% ============================================================
% STAGE 5: Label Preservation
% ============================================================
\node[mainblock,
  right=0.30cm of valid,
  fill=gray!8, draw=gray!50,
  minimum height=1.2cm
] (labels) {Label\\Preservation};

\node[subblock, below=0.08cm of labels] (labels_sub)
  {VeReMi kinematic\\label kept separate\\never merged};

% ============================================================
% STAGE 6: STBV-Bench v1
% [LAYOUT] Bench block kept at original 3.0cm width (unchanged)
% ============================================================
\node[mainblock,
  right=0.30cm of labels,
  fill=benchpurple!10, draw=benchpurple, text=benchpurple,
  minimum height=1.8cm, minimum width=3.0cm
] (bench) {};

\node[font=\scriptsize\bfseries, text=benchpurple, align=center,
      at=(bench.center), yshift=0.60cm]
  {STBV-Bench v1};
\node[font=\tiny, text=benchpurple!80, align=center,
      at=(bench.center), yshift=0.25cm]
  {\textbf{n = 10,000 samples}};
\node[font=\tiny, text=benchpurple!80, align=center,
      at=(bench.center), yshift=-0.05cm]
  {70.07\% malicious};
\node[font=\tiny, text=benchpurple!80, align=center,
      at=(bench.center), yshift=-0.35cm]
  {29.93\% benign};
\node[font=\tiny, text=benchpurple!80, align=center,
      at=(bench.center), yshift=-0.65cm]
  {seed = fixed};

% ============================================================
% STAGE 7: Evaluation Pipeline
% ============================================================
\node[mainblock,
  right=0.30cm of bench,
  fill=evalcyan!10, draw=evalcyan, text=evalcyan,
  minimum height=1.2cm
] (eval) {STBV\\Architecture};

\node[subblock, below=0.08cm of eval] (eval_sub)
  {B1+B2+B3\\+Fusion\\per-message\\independent};

% ============================================================
% ARROWS
% ============================================================
\draw[arrow, veremiblue!60] (veremi.east)  -- (extract.west);
\draw[arrow]                (extract.east) -- (engine.west);
\draw[arrow, engineorange!60] (engine.east) -- (valid.west);
\draw[arrow]                (valid.east)   -- (labels.west);
\draw[arrow, benchpurple!60] (labels.east) -- (bench.west);
\draw[arrow, evalcyan!60]   (bench.east)   -- (eval.west);

% ============================================================
% KINEMATIC BYPASS ARROW (kinematics pass through unmodified)
% [LAYOUT] Label moved from pos=0.25 to pos=0.5 for better
%          centering under the dashed path
% ============================================================
\draw[
  -{Stealth[length=4pt]}, dashed, gray!40, thin
] (extract_sub.south) -- ++(0,-0.4)
  -| node[pos=0.5, below, font=\tiny, text=black]
     {kinematics unmodified (pos, speed, heading, timestamp)}
  (labels_sub.south);

% ============================================================
% BOTTOM STATS BAR
% ============================================================
\node[
  font=\tiny, text=black, align=center,
  below=1.2cm of engine
] {Pool: 221,125 records \quad$\xrightarrow{\text{sample}}$\quad
   n=10,000 \quad$\xrightarrow{\text{evaluate}}$\quad
   RQ1--RQ6};

\end{tikzpicture}

%     \caption{STBV-Bench generation pipeline. Real VeReMi Extension
%     trajectories supply kinematics; the seeded semantic transformation
%     engine populates free-text scene-context fields across 20 attack
%     families plus a benign control. Kinematic ground truth is preserved
%     separately and never merged with semantic labels.}
%     \label{fig:stbvbench_pipeline}
%   \end{figure*}
% ============================================================

\begin{tikzpicture}[
  arrow/.style={-{Stealth[length=5pt]}, thick, gray!55},
  % [LAYOUT] Stage 1 (Dataset) and Stage 6 (STBV-Bench) blocks kept at
  %          their original widths (2.8cm and 3.0cm respectively).
  %          Intermediate stages reduced from 2.8cm to 2.3cm to recover
  %          horizontal space without altering the visual hierarchy.
  %          The Engine remains the widest and tallest block.
  mainblock/.style={
    rectangle, rounded corners=3pt,
    minimum width=1.85cm, minimum height=1.0cm,
    font=\scriptsize\bfseries, align=center,
    draw, thick
  },
  subblock/.style={
    rectangle, rounded corners=2pt,
    minimum width=1.8cm, minimum height=0.48cm,
    font=\tiny, align=center, text width=1.9cm,
    draw=gray!50, fill=gray!8, thin
  },
  statnode/.style={
    font=\tiny\bfseries, align=center, text width=2.4cm
  },
  % [LAYOUT] Reduced inter-block spacing from 0.55cm to 0.35cm.
  node distance=0.5cm and 0.35cm
]

\definecolor{veremiblue}{RGB}{60,100,180}
\definecolor{engineorange}{RGB}{200,100,20}
\definecolor{validgreen}{RGB}{40,140,60}
\definecolor{benchpurple}{RGB}{120,40,160}
\definecolor{evalcyan}{RGB}{20,120,140}

% ============================================================
% STAGE 1: VeReMi Source
% [LAYOUT] Dataset block kept at original 2.8cm width (unchanged)
% ============================================================
\node[mainblock,
  fill=veremiblue!12, draw=veremiblue, text=veremiblue,
  minimum width=2.8cm, minimum height=1.2cm
] (veremi) {VeReMi Extension\\Dataset};

\node[statnode, below=0.10cm of veremi] (veremi_stat)
  {\textcolor{veremiblue}{221,125 records}\\
   \textcolor{black}{360 vehicles}\\
   \textcolor{black}{180 attackers}};

% ============================================================
% STAGE 2: Field Extraction
% ============================================================
\node[mainblock,
  right=0.30cm of veremi,
  fill=gray!8, draw=gray!50,
  minimum height=1.2cm
] (extract) {Field\\Extraction};

\node[subblock, below=0.08cm of extract] (extract_sub)
  {pos, speed, heading,\\timestamp\\(unmodified)};

% ============================================================
% STAGE 3: Semantic Transformation Engine
% [LAYOUT] Engine kept as largest/widest block at 3.2cm wide,
%          2.2cm tall, to preserve its focal role in the figure
% ============================================================
\node[mainblock,
  right=0.30cm of extract,
  fill=engineorange!10, draw=engineorange, text=engineorange,
  minimum height=3.2cm, minimum width=3.8cm
] (engine) {};

% Explicitly position every title line to guarantee no overlaps
\node[font=\scriptsize\bfseries, text=engineorange, align=center,
      at=(engine.center), yshift=1.0cm]
  {Semantic};
\node[font=\scriptsize\bfseries, text=engineorange, align=center,
      at=(engine.center), yshift=0.65cm]
  {Transformation};
\node[font=\scriptsize\bfseries, text=engineorange, align=center,
      at=(engine.center), yshift=0.30cm]
  {Engine};

% Sub-labels inside engine — positioned explicitly below the title
\node[font=\tiny, text=engineorange!80, align=center,
      at=(engine.center), yshift=-0.25cm]
  {seed = fixed};
\node[font=\tiny, text=engineorange!80, align=center,
      at=(engine.center), yshift=-0.65cm]
  {rule-based, templated};
\node[font=\tiny, text=engineorange!80, align=center,
      at=(engine.center), yshift=-1.05cm]
  {no LLM, no paraphrase};

% Attack families annotation (above engine)
\node[
  font=\tiny\bfseries, text=engineorange,
  align=center, text width=3.2cm,
  above=0.15cm of engine
] {20 attack families + 1 benign control = 21 transformation rules};

% ============================================================
% STAGE 4: Structural Validation
% ============================================================
\node[mainblock,
  right=0.30cm of engine,
  fill=validgreen!10, draw=validgreen, text=validgreen,
  minimum height=1.2cm
] (valid) {Structural\\Validation};

\node[subblock, below=0.08cm of valid] (valid_sub)
  {schema check\\abort on failure};

% ============================================================
% STAGE 5: Label Preservation
% ============================================================
\node[mainblock,
  right=0.30cm of valid,
  fill=gray!8, draw=gray!50,
  minimum height=1.2cm
] (labels) {Label\\Preservation};

\node[subblock, below=0.08cm of labels] (labels_sub)
  {VeReMi kinematic\\label kept separate\\never merged};

% ============================================================
% STAGE 6: STBV-Bench v1
% [LAYOUT] Bench block kept at original 3.0cm width (unchanged)
% ============================================================
\node[mainblock,
  right=0.30cm of labels,
  fill=benchpurple!10, draw=benchpurple, text=benchpurple,
  minimum height=1.8cm, minimum width=3.0cm
] (bench) {};

\node[font=\scriptsize\bfseries, text=benchpurple, align=center,
      at=(bench.center), yshift=0.60cm]
  {STBV-Bench v1};
\node[font=\tiny, text=benchpurple!80, align=center,
      at=(bench.center), yshift=0.25cm]
  {\textbf{n = 10,000 samples}};
\node[font=\tiny, text=benchpurple!80, align=center,
      at=(bench.center), yshift=-0.05cm]
  {70.07\% malicious};
\node[font=\tiny, text=benchpurple!80, align=center,
      at=(bench.center), yshift=-0.35cm]
  {29.93\% benign};
\node[font=\tiny, text=benchpurple!80, align=center,
      at=(bench.center), yshift=-0.65cm]
  {seed = fixed};

% ============================================================
% STAGE 7: Evaluation Pipeline
% ============================================================
\node[mainblock,
  right=0.30cm of bench,
  fill=evalcyan!10, draw=evalcyan, text=evalcyan,
  minimum height=1.2cm
] (eval) {STBV\\Architecture};

\node[subblock, below=0.08cm of eval] (eval_sub)
  {B1+B2+B3\\+Fusion\\per-message\\independent};

% ============================================================
% ARROWS
% ============================================================
\draw[arrow, veremiblue!60] (veremi.east)  -- (extract.west);
\draw[arrow]                (extract.east) -- (engine.west);
\draw[arrow, engineorange!60] (engine.east) -- (valid.west);
\draw[arrow]                (valid.east)   -- (labels.west);
\draw[arrow, benchpurple!60] (labels.east) -- (bench.west);
\draw[arrow, evalcyan!60]   (bench.east)   -- (eval.west);

% ============================================================
% KINEMATIC BYPASS ARROW (kinematics pass through unmodified)
% [LAYOUT] Label moved from pos=0.25 to pos=0.5 for better
%          centering under the dashed path
% ============================================================
\draw[
  -{Stealth[length=4pt]}, dashed, gray!40, thin
] (extract_sub.south) -- ++(0,-0.4)
  -| node[pos=0.5, below, font=\tiny, text=black]
     {kinematics unmodified (pos, speed, heading, timestamp)}
  (labels_sub.south);

% ============================================================
% BOTTOM STATS BAR
% ============================================================
\node[
  font=\tiny, text=black, align=center,
  below=1.2cm of engine
] {Pool: 221,125 records \quad$\xrightarrow{\text{sample}}$\quad
   n=10,000 \quad$\xrightarrow{\text{evaluate}}$\quad
   RQ1--RQ6};

\end{tikzpicture}

%     \caption{STBV-Bench generation pipeline. Real VeReMi Extension
%     trajectories supply kinematics; the seeded semantic transformation
%     engine populates free-text scene-context fields across 20 attack
%     families plus a benign control. Kinematic ground truth is preserved
%     separately and never merged with semantic labels.}
%     \label{fig:stbvbench_pipeline}
%   \end{figure*}
% ============================================================

\begin{tikzpicture}[
  arrow/.style={-{Stealth[length=5pt]}, thick, gray!55},
  % [LAYOUT] Stage 1 (Dataset) and Stage 6 (STBV-Bench) blocks kept at
  %          their original widths (2.8cm and 3.0cm respectively).
  %          Intermediate stages reduced from 2.8cm to 2.3cm to recover
  %          horizontal space without altering the visual hierarchy.
  %          The Engine remains the widest and tallest block.
  mainblock/.style={
    rectangle, rounded corners=3pt,
    minimum width=1.85cm, minimum height=1.0cm,
    font=\scriptsize\bfseries, align=center,
    draw, thick
  },
  subblock/.style={
    rectangle, rounded corners=2pt,
    minimum width=1.8cm, minimum height=0.48cm,
    font=\tiny, align=center, text width=1.9cm,
    draw=gray!50, fill=gray!8, thin
  },
  statnode/.style={
    font=\tiny\bfseries, align=center, text width=2.4cm
  },
  % [LAYOUT] Reduced inter-block spacing from 0.55cm to 0.35cm.
  node distance=0.5cm and 0.35cm
]

\definecolor{veremiblue}{RGB}{60,100,180}
\definecolor{engineorange}{RGB}{200,100,20}
\definecolor{validgreen}{RGB}{40,140,60}
\definecolor{benchpurple}{RGB}{120,40,160}
\definecolor{evalcyan}{RGB}{20,120,140}

% ============================================================
% STAGE 1: VeReMi Source
% [LAYOUT] Dataset block kept at original 2.8cm width (unchanged)
% ============================================================
\node[mainblock,
  fill=veremiblue!12, draw=veremiblue, text=veremiblue,
  minimum width=2.8cm, minimum height=1.2cm
] (veremi) {VeReMi Extension\\Dataset};

\node[statnode, below=0.10cm of veremi] (veremi_stat)
  {\textcolor{veremiblue}{221,125 records}\\
   \textcolor{black}{360 vehicles}\\
   \textcolor{black}{180 attackers}};

% ============================================================
% STAGE 2: Field Extraction
% ============================================================
\node[mainblock,
  right=0.30cm of veremi,
  fill=gray!8, draw=gray!50,
  minimum height=1.2cm
] (extract) {Field\\Extraction};

\node[subblock, below=0.08cm of extract] (extract_sub)
  {pos, speed, heading,\\timestamp\\(unmodified)};

% ============================================================
% STAGE 3: Semantic Transformation Engine
% [LAYOUT] Engine kept as largest/widest block at 3.2cm wide,
%          2.2cm tall, to preserve its focal role in the figure
% ============================================================
\node[mainblock,
  right=0.30cm of extract,
  fill=engineorange!10, draw=engineorange, text=engineorange,
  minimum height=3.2cm, minimum width=3.8cm
] (engine) {};

% Explicitly position every title line to guarantee no overlaps
\node[font=\scriptsize\bfseries, text=engineorange, align=center,
      at=(engine.center), yshift=1.0cm]
  {Semantic};
\node[font=\scriptsize\bfseries, text=engineorange, align=center,
      at=(engine.center), yshift=0.65cm]
  {Transformation};
\node[font=\scriptsize\bfseries, text=engineorange, align=center,
      at=(engine.center), yshift=0.30cm]
  {Engine};

% Sub-labels inside engine — positioned explicitly below the title
\node[font=\tiny, text=engineorange!80, align=center,
      at=(engine.center), yshift=-0.25cm]
  {seed = fixed};
\node[font=\tiny, text=engineorange!80, align=center,
      at=(engine.center), yshift=-0.65cm]
  {rule-based, templated};
\node[font=\tiny, text=engineorange!80, align=center,
      at=(engine.center), yshift=-1.05cm]
  {no LLM, no paraphrase};

% Attack families annotation (above engine)
\node[
  font=\tiny\bfseries, text=engineorange,
  align=center, text width=3.2cm,
  above=0.15cm of engine
] {20 attack families + 1 benign control = 21 transformation rules};

% ============================================================
% STAGE 4: Structural Validation
% ============================================================
\node[mainblock,
  right=0.30cm of engine,
  fill=validgreen!10, draw=validgreen, text=validgreen,
  minimum height=1.2cm
] (valid) {Structural\\Validation};

\node[subblock, below=0.08cm of valid] (valid_sub)
  {schema check\\abort on failure};

% ============================================================
% STAGE 5: Label Preservation
% ============================================================
\node[mainblock,
  right=0.30cm of valid,
  fill=gray!8, draw=gray!50,
  minimum height=1.2cm
] (labels) {Label\\Preservation};

\node[subblock, below=0.08cm of labels] (labels_sub)
  {VeReMi kinematic\\label kept separate\\never merged};

% ============================================================
% STAGE 6: STBV-Bench v1
% [LAYOUT] Bench block kept at original 3.0cm width (unchanged)
% ============================================================
\node[mainblock,
  right=0.30cm of labels,
  fill=benchpurple!10, draw=benchpurple, text=benchpurple,
  minimum height=1.8cm, minimum width=3.0cm
] (bench) {};

\node[font=\scriptsize\bfseries, text=benchpurple, align=center,
      at=(bench.center), yshift=0.60cm]
  {STBV-Bench v1};
\node[font=\tiny, text=benchpurple!80, align=center,
      at=(bench.center), yshift=0.25cm]
  {\textbf{n = 10,000 samples}};
\node[font=\tiny, text=benchpurple!80, align=center,
      at=(bench.center), yshift=-0.05cm]
  {70.07\% malicious};
\node[font=\tiny, text=benchpurple!80, align=center,
      at=(bench.center), yshift=-0.35cm]
  {29.93\% benign};
\node[font=\tiny, text=benchpurple!80, align=center,
      at=(bench.center), yshift=-0.65cm]
  {seed = fixed};

% ============================================================
% STAGE 7: Evaluation Pipeline
% ============================================================
\node[mainblock,
  right=0.30cm of bench,
  fill=evalcyan!10, draw=evalcyan, text=evalcyan,
  minimum height=1.2cm
] (eval) {STBV\\Architecture};

\node[subblock, below=0.08cm of eval] (eval_sub)
  {B1+B2+B3\\+Fusion\\per-message\\independent};

% ============================================================
% ARROWS
% ============================================================
\draw[arrow, veremiblue!60] (veremi.east)  -- (extract.west);
\draw[arrow]                (extract.east) -- (engine.west);
\draw[arrow, engineorange!60] (engine.east) -- (valid.west);
\draw[arrow]                (valid.east)   -- (labels.west);
\draw[arrow, benchpurple!60] (labels.east) -- (bench.west);
\draw[arrow, evalcyan!60]   (bench.east)   -- (eval.west);

% ============================================================
% KINEMATIC BYPASS ARROW (kinematics pass through unmodified)
% [LAYOUT] Label moved from pos=0.25 to pos=0.5 for better
%          centering under the dashed path
% ============================================================
\draw[
  -{Stealth[length=4pt]}, dashed, gray!40, thin
] (extract_sub.south) -- ++(0,-0.4)
  -| node[pos=0.5, below, font=\tiny, text=black]
     {kinematics unmodified (pos, speed, heading, timestamp)}
  (labels_sub.south);

% ============================================================
% BOTTOM STATS BAR
% ============================================================
\node[
  font=\tiny, text=black, align=center,
  below=1.2cm of engine
] {Pool: 221,125 records \quad$\xrightarrow{\text{sample}}$\quad
   n=10,000 \quad$\xrightarrow{\text{evaluate}}$\quad
   RQ1--RQ6};

\end{tikzpicture}

%     \caption{STBV-Bench generation pipeline. Real VeReMi Extension
%     trajectories supply kinematics; the seeded semantic transformation
%     engine populates free-text scene-context fields across 20 attack
%     families plus a benign control. Kinematic ground truth is preserved
%     separately and never merged with semantic labels.}
%     \label{fig:stbvbench_pipeline}
%   \end{figure*}
% ============================================================

\begin{tikzpicture}[
  arrow/.style={-{Stealth[length=5pt]}, thick, gray!55},
  % [LAYOUT] Stage 1 (Dataset) and Stage 6 (STBV-Bench) blocks kept at
  %          their original widths (2.8cm and 3.0cm respectively).
  %          Intermediate stages reduced from 2.8cm to 2.3cm to recover
  %          horizontal space without altering the visual hierarchy.
  %          The Engine remains the widest and tallest block.
  mainblock/.style={
    rectangle, rounded corners=3pt,
    minimum width=1.85cm, minimum height=1.0cm,
    font=\scriptsize\bfseries, align=center,
    draw, thick
  },
  subblock/.style={
    rectangle, rounded corners=2pt,
    minimum width=1.8cm, minimum height=0.48cm,
    font=\tiny, align=center, text width=1.9cm,
    draw=gray!50, fill=gray!8, thin
  },
  statnode/.style={
    font=\tiny\bfseries, align=center, text width=2.4cm
  },
  % [LAYOUT] Reduced inter-block spacing from 0.55cm to 0.35cm.
  node distance=0.5cm and 0.35cm
]

\definecolor{veremiblue}{RGB}{60,100,180}
\definecolor{engineorange}{RGB}{200,100,20}
\definecolor{validgreen}{RGB}{40,140,60}
\definecolor{benchpurple}{RGB}{120,40,160}
\definecolor{evalcyan}{RGB}{20,120,140}

% ============================================================
% STAGE 1: VeReMi Source
% [LAYOUT] Dataset block kept at original 2.8cm width (unchanged)
% ============================================================
\node[mainblock,
  fill=veremiblue!12, draw=veremiblue, text=veremiblue,
  minimum width=2.8cm, minimum height=1.2cm
] (veremi) {VeReMi Extension\\Dataset};

\node[statnode, below=0.10cm of veremi] (veremi_stat)
  {\textcolor{veremiblue}{221,125 records}\\
   \textcolor{black}{360 vehicles}\\
   \textcolor{black}{180 attackers}};

% ============================================================
% STAGE 2: Field Extraction
% ============================================================
\node[mainblock,
  right=0.30cm of veremi,
  fill=gray!8, draw=gray!50,
  minimum height=1.2cm
] (extract) {Field\\Extraction};

\node[subblock, below=0.08cm of extract] (extract_sub)
  {pos, speed, heading,\\timestamp\\(unmodified)};

% ============================================================
% STAGE 3: Semantic Transformation Engine
% [LAYOUT] Engine kept as largest/widest block at 3.2cm wide,
%          2.2cm tall, to preserve its focal role in the figure
% ============================================================
\node[mainblock,
  right=0.30cm of extract,
  fill=engineorange!10, draw=engineorange, text=engineorange,
  minimum height=3.2cm, minimum width=3.8cm
] (engine) {};

% Explicitly position every title line to guarantee no overlaps
\node[font=\scriptsize\bfseries, text=engineorange, align=center,
      at=(engine.center), yshift=1.0cm]
  {Semantic};
\node[font=\scriptsize\bfseries, text=engineorange, align=center,
      at=(engine.center), yshift=0.65cm]
  {Transformation};
\node[font=\scriptsize\bfseries, text=engineorange, align=center,
      at=(engine.center), yshift=0.30cm]
  {Engine};

% Sub-labels inside engine — positioned explicitly below the title
\node[font=\tiny, text=engineorange!80, align=center,
      at=(engine.center), yshift=-0.25cm]
  {seed = fixed};
\node[font=\tiny, text=engineorange!80, align=center,
      at=(engine.center), yshift=-0.65cm]
  {rule-based, templated};
\node[font=\tiny, text=engineorange!80, align=center,
      at=(engine.center), yshift=-1.05cm]
  {no LLM, no paraphrase};

% Attack families annotation (above engine)
\node[
  font=\tiny\bfseries, text=engineorange,
  align=center, text width=3.2cm,
  above=0.15cm of engine
] {20 attack families + 1 benign control = 21 transformation rules};

% ============================================================
% STAGE 4: Structural Validation
% ============================================================
\node[mainblock,
  right=0.30cm of engine,
  fill=validgreen!10, draw=validgreen, text=validgreen,
  minimum height=1.2cm
] (valid) {Structural\\Validation};

\node[subblock, below=0.08cm of valid] (valid_sub)
  {schema check\\abort on failure};

% ============================================================
% STAGE 5: Label Preservation
% ============================================================
\node[mainblock,
  right=0.30cm of valid,
  fill=gray!8, draw=gray!50,
  minimum height=1.2cm
] (labels) {Label\\Preservation};

\node[subblock, below=0.08cm of labels] (labels_sub)
  {VeReMi kinematic\\label kept separate\\never merged};

% ============================================================
% STAGE 6: STBV-Bench v1
% [LAYOUT] Bench block kept at original 3.0cm width (unchanged)
% ============================================================
\node[mainblock,
  right=0.30cm of labels,
  fill=benchpurple!10, draw=benchpurple, text=benchpurple,
  minimum height=1.8cm, minimum width=3.0cm
] (bench) {};

\node[font=\scriptsize\bfseries, text=benchpurple, align=center,
      at=(bench.center), yshift=0.60cm]
  {STBV-Bench v1};
\node[font=\tiny, text=benchpurple!80, align=center,
      at=(bench.center), yshift=0.25cm]
  {\textbf{n = 10,000 samples}};
\node[font=\tiny, text=benchpurple!80, align=center,
      at=(bench.center), yshift=-0.05cm]
  {70.07\% malicious};
\node[font=\tiny, text=benchpurple!80, align=center,
      at=(bench.center), yshift=-0.35cm]
  {29.93\% benign};
\node[font=\tiny, text=benchpurple!80, align=center,
      at=(bench.center), yshift=-0.65cm]
  {seed = fixed};

% ============================================================
% STAGE 7: Evaluation Pipeline
% ============================================================
\node[mainblock,
  right=0.30cm of bench,
  fill=evalcyan!10, draw=evalcyan, text=evalcyan,
  minimum height=1.2cm
] (eval) {STBV\\Architecture};

\node[subblock, below=0.08cm of eval] (eval_sub)
  {B1+B2+B3\\+Fusion\\per-message\\independent};

% ============================================================
% ARROWS
% ============================================================
\draw[arrow, veremiblue!60] (veremi.east)  -- (extract.west);
\draw[arrow]                (extract.east) -- (engine.west);
\draw[arrow, engineorange!60] (engine.east) -- (valid.west);
\draw[arrow]                (valid.east)   -- (labels.west);
\draw[arrow, benchpurple!60] (labels.east) -- (bench.west);
\draw[arrow, evalcyan!60]   (bench.east)   -- (eval.west);

% ============================================================
% KINEMATIC BYPASS ARROW (kinematics pass through unmodified)
% [LAYOUT] Label moved from pos=0.25 to pos=0.5 for better
%          centering under the dashed path
% ============================================================
\draw[
  -{Stealth[length=4pt]}, dashed, gray!40, thin
] (extract_sub.south) -- ++(0,-0.4)
  -| node[pos=0.5, below, font=\tiny, text=black]
     {kinematics unmodified (pos, speed, heading, timestamp)}
  (labels_sub.south);

% ============================================================
% BOTTOM STATS BAR
% ============================================================
\node[
  font=\tiny, text=black, align=center,
  below=1.2cm of engine
] {Pool: 221,125 records \quad$\xrightarrow{\text{sample}}$\quad
   n=10,000 \quad$\xrightarrow{\text{evaluate}}$\quad
   RQ1--RQ6};

\end{tikzpicture}
